\subsection*{Purpose of this Document}
This document outlines the technical principles, architecture, and integration strategy of ScienceEcosystem. It is intended to guide implementation, audits, and long-term maintenance while remaining adaptable to future standards.

\section*{1. Architectural Philosophy}

ScienceEcosystem is designed as a connective infrastructure layer, not a monolithic platform. Its role is to link, contextualise, and preserve relationships between existing research components.

Key design goals:
\begin{itemize}
    \item Modularity
    \item Interoperability
    \item Durability
    \item Auditability
\end{itemize}

\section*{2. Core System Components}

\subsection*{2.1 Metadata Layer}
\begin{itemize}
    \item Persistent identifiers for papers, datasets, code, authors, and institutions
    \item Normalised metadata schemas
    \item Versioned provenance records
\end{itemize}

\subsection*{2.2 Relationship Graph}
\begin{itemize}
    \item Explicit links between publications, data, code, and outputs
    \item Support for temporal and versioned relationships
    \item Provenance trails for verification
\end{itemize}

\subsection*{2.3 Storage Strategy}
ScienceEcosystem does not duplicate primary storage unnecessarily.
\begin{itemize}
    \item External repositories remain authoritative
    \item Archival snapshots referenced where required
    \item Long-term preservation planning integrated
\end{itemize}

\subsection*{2.4 Search and Indexing}
\begin{itemize}
    \item Unified search across metadata and relationships
    \item Transparent ranking criteria
    \item No algorithmic promotion or suppression
\end{itemize}

\section*{3. Integration Strategy}

ScienceEcosystem integrates with existing infrastructure:
\begin{itemize}
    \item Source control systems for code
    \item Archival repositories for datasets
    \item Identifier services for authors and institutions
    \item Development environments for reproducible workflows
\end{itemize}

ScienceEcosystem acts as the research object and provenance layer binding these components.

\section*{4. Reproducibility and Verification}

The platform supports reproducibility by:
\begin{itemize}
    \item Linking code and data at version level
    \item Recording computational environments where possible
    \item Allowing re-execution and verification workflows
\end{itemize}

Verification is optional but structurally enabled.

\section*{5. Security and Privacy}

\begin{itemize}
    \item Compliance with international data protection regulations
    \item Minimal collection of personal data
    \item Regular security audits
\end{itemize}

No tracking or behavioural profiling is performed.

\section*{6. Open Source and Forkability}

Core components are open-source where feasible. Documentation and deployment scripts are public to:
\begin{itemize}
    \item Enable independent audits
    \item Prevent platform capture
    \item Ensure continuity beyond any single operator
\end{itemize}

\section*{7. Scalability and Maintenance}

Scalability is achieved through:
\begin{itemize}
    \item Stateless services
    \item Horizontal scaling
    \item Conservative technology choices
\end{itemize}

Maintenance prioritises reliability over feature velocity.

\section*{8. Technical Governance}

Major architectural changes require:
\begin{itemize}
    \item Public technical proposals
    \item Review by qualified contributors
    \item Alignment with governance principles
\end{itemize}

\section*{9. Commitment}

ScienceEcosystem’s technical design exists to make science verifiable, not optimised for attention or speed. Stability and trust outweigh novelty.


